\documentclass[12pt,a4paper]{article}

% Paquetes básicos
\usepackage[table,x11names]{xcolor}
\usepackage{alltt}
\usepackage[utf8]{inputenc}
\usepackage{mathtools, amssymb, amsthm}
\usepackage{changepage}
\usepackage{geometry}
\usepackage[colorlinks=true]{hyperref}
\usepackage{enumitem}
\usepackage{etoolbox}
\usepackage{graphicx}
\usepackage{setspace}
\usepackage{tcolorbox}
\tcbuselibrary{skins, breakable}
\usepackage{titlesec}
\usepackage{float}

\usepackage{listings}

\lstset{
  basicstyle=\ttfamily,
  columns=fullflexible
}

% Margins
%\geometry{left=3cm,right=3cm,top=2.5cm,bottom=2.5cm}

% Custom operators
\newcommand{\card}{\operatorname{card}}
\newcommand{\muae}{\overset{\mu-a.e.}{=}}

% Implication box setup
\tcbset{Implication-number/.style={
  enhanced,
  boxsep=2pt,
  colback=white,
  frame hidden,
  sharp corners,
  left=2pt, right=2pt, top=1pt, bottom=1pt,
  underlay={
    \draw[line width=0.5pt] (frame.south west) -- ([xshift=-133mm]frame.south east); % línea horizontal
    \draw[line width=0.5pt] ([xshift=-133mm]frame.north east) -- ([xshift=-133mm]frame.south east); % línea vertical
  }
}}

\tcbset{Implication-number-ds/.style={
  enhanced,
  boxsep=2pt,
  colback=white,
  frame hidden,
  sharp corners,
  left=2pt, right=2pt, top=1pt, bottom=1pt,
  underlay={
    \draw[line width=0.5pt] ([yshift=5mm]frame.south west) -- ([xshift=-133mm, yshift=5mm]frame.south east); % línea horizontal
    \draw[line width=0.5pt] ([xshift=-133mm, yshift=-3mm]frame.north east) -- ([xshift=-133mm, yshift=5mm]frame.south east); % línea vertical
  }
}}

\tcbset{Subset-contingency/.style={
  enhanced,
  boxsep=2pt,
  colback=white,
  frame hidden,
  sharp corners,
  left=2pt, right=2pt, top=1pt, bottom=1pt,
  underlay={
    \draw[line width=0.5pt] (frame.south west) -- ([xshift=-140mm]frame.south east); % línea horizontal
    \draw[line width=0.5pt] ([xshift=-140mm]frame.north east) -- ([xshift=-140mm]frame.south east); % línea vertical
  }
}}

\tcbset{Indent-subset-contingency/.style={
  enhanced,
  boxsep=2pt,
  colback=white,
  frame hidden,
  sharp corners,
  left=2pt, right=2pt, top=1pt, bottom=1pt,
  underlay={
    \draw[line width=0.5pt] (frame.south west) -- ([xshift=-140mm+0.07\textwidth]frame.south east); % línea horizontal
    \draw[line width=0.5pt] ([xshift=-140mm+0.07\textwidth]frame.north east) -- ([xshift=-140mm+0.07\textwidth]frame.south east); % línea vertical
  }
}}

% Useful commands
\renewcommand{\contentsname}{Contenidos}

\newcommand{\R}{\mathbb{R}}
\newcommand{\N}{\mathbb{N}}
\newcommand{\Z}{\mathbb{Z}}
\newcommand{\Q}{\mathbb{Q}}
\newcommand{\C}{\mathbb{C}}

\newcommand{\smallcup}{\mathop{\cup}\limits}
\newcommand{\smallcap}{\mathop{\cap}\limits}
\newcommand{\smallsum}{\mathop{\sum}\limits}
\newcommand{\smallprod}{\mathop{\prod}\limits}

\newcommand{\linf}[1]{\displaystyle{\mathop{\underline{\lim}}_{#1}}}
\newcommand{\mlim}[1]{\displaystyle{\lim_{#1}}}

%Integral de Lebesgue con patas
\newcommand{\lbint}{\mathop{\int_{\!\!\!\!\!|\!}^{\!|\!}}}

%Espacios \mathcal{L}_p
\newcommand{\elep}[2]{\mathcal{L}_{#1}(#2)}

% ----- Custom counters and counter commands -----
% Custom counter hierarchy
\newcounter{unit}[section]
\newcounter{chapter}[unit]
\makeatletter
\@addtoreset{subsubsection}{chapter}
\makeatother

\renewcommand{\theunit}{\arabic{unit}}
\renewcommand{\thechapter}{\arabic{chapter}}
\renewcommand{\thesubsubsection}{\arabic{chapter}.\arabic{subsubsection}.}

% Custom content hierarchy behavior
\newcommand{\chapter}[1]{
    \refstepcounter{chapter}
    \subsection*{\Large{\S \thechapter. #1}}
    \addcontentsline{toc}{subsection}{\thechapter. #1}
}
\newcommand{\unit}[1]{
    \refstepcounter{unit}
    \section*{\Huge{\Roman{unit} #1}}
    \addcontentsline{toc}{section}{\Roman{unit} #1}
}

\newcommand{\result}[1]{%
  \subsubsection{#1}%
  \label{result:\thesubsubsection}
}
  
\titleformat{\subsubsection}
    {\normalfont\large\bfseries} % mismo tamaño que \subsection
    {\thesubsubsection}{1em}{}
  
%---- Custom proof commands-----
\newcommand{\dem}{
    \noindent \underline{\textbf{Demostración:}}
}
\newcommand{\nota}{
    \noindent \underline{\textbf{Nota:}}
}
% ----------------------------------------
\hbadness=10000
\vbadness=10000
\hfuzz=100pt
\vfuzz=100pt
% ---------------------------------------
\title{Algoritmia}
\author{Práctica 3}
\date{21/2/2026}


\begin{document}

\maketitle
\hypersetup{linkcolor=black}
\vspace{4mm}
\tableofcontents
\hypersetup{linkcolor=Ivory4}
\newpage

\chapter{Divide y vencerás por sustracción}
\result{Complejidades teóricas de los algoritmos}
Usaremos el Teorema Maestro para calcular la complejidad temporal para cada método recursivo de esta práctica. Siguiendo la notación usual,
\begin{itemize}
  \item \textbf{Sustracción 1:} 
  \begin{flalign*}
      \begin{cases}
        k = 0\\
        a = 1\\
        b = 1
      \end{cases} \Rightarrow T_{S_1}(n) = O(n) &&
  \end{flalign*}
  \item \textbf{Sustracción 2:} 
  \begin{flalign*}
      \begin{cases}
        k = 1\\
        a = 1\\
        b = 1
      \end{cases} \Rightarrow T_{S_2}(n) = O(n^2) &&
  \end{flalign*}
  \item \textbf{Sustracción 3:} 
  \begin{flalign*}
      \begin{cases}
        k = 0\\
        a = 2\\
        b = 1
      \end{cases} \Rightarrow T_{S_3}(n) = O\left(2^n\right) &&
  \end{flalign*}
  \item \textbf{Sustracción 4:} 
  \begin{flalign*}
      \begin{cases}
        k = 2\\
        a = 1\\
        b = 1
      \end{cases} \Rightarrow T_{S_4}(n) = O(n^3) &&
  \end{flalign*}
  \item \textbf{Sustracción 5:} 
  \begin{flalign*}
      \begin{cases}
        k = 0\\
        a = 3\\
        b = 2
      \end{cases} \Rightarrow T_{S_5}(n) = O\left(3^\frac{n}{2}\right) &&
  \end{flalign*}
\end{itemize}

\newpage
\result{Valores de \texorpdfstring{$n$}{n} para los que se aborta la ejecución}
\hspace{3mm} Ejecutando las clases \texttt{Sustraccion1} y \texttt{Sustraccion2} (con una sola repetición en cada caso), podemos observar que se aborta la ejecución:

\begin{itemize}
  \item \textbf{Sustracción 1:} \\[-4ex]
  \begin{alltt}
 n=1000 TIEMPO=0  nVeces=1
 n=2000 TIEMPO=0  nVeces=1
 n=4000 TIEMPO=0  nVeces=1
 n=8000 TIEMPO=1  nVeces=1
 \textcolor{Firebrick2}{Exception in thread "main" java.lang.StackOverflowError}
\end{alltt}

\item \textbf{Sustracción 2:} \\[-4ex]
\begin{alltt}
 n=1000 TIEMPO=6  nVeces=1
 n=2000 TIEMPO=27  nVeces=1
 n=4000 TIEMPO=104  nVeces=1
 n=8000 TIEMPO=400  nVeces=1
 \textcolor{Firebrick2}{Exception in thread "main" java.lang.StackOverflowError}
\end{alltt}
\end{itemize}

Para los dos métodos, la altura del árbol de llamadas es tal que se llega a un desbordamiento de pila. En ambos casos, cada llamada al método correspondiente con tamaño $n$ induce otra de tamaño $n-1$ hasta que el problema tenga un tamaño fijo (caso base). Así, la altura del arbol de llamadas es lineal: $O(n)$.

\vspace{6mm}
\result{Final de \texttt{sustraccion3} para tamaño \texorpdfstring{$n=80$}{n=80}}
\hspace{3mm} Basándonos en los datos empíricos medidos para este método, para el caso $n=29$ tenemos un tiempo de ejecución de $22.057ms$. Comparando ahora con la complejidad teórica:
$$t_2 = t_1 \cdot \frac{f(n_2)}{f(n_1)} = 22.057ms \cdot \frac{2^{80}}{2^{29}}\cdot\frac{1h}{3,6 \cdot 10^6}ms \cdot \frac{1\text{año}}{24*365h}= 1,5750 \cdot 10^9 \text{años}$$

\newpage
\result{Toma de tiempos para \texttt{sustraccion4}}
\hspace{3mm} Tal y como hemos mencionado en el \hyperref[result:1.1.]{apartado 1.1}, para lograr la complejidad teórica de $O(n^3)$, hemos empleado recursividad por sustracción haciendo uso de los siguientes parámetros:
\begin{align*}
    k = 2 && a = 1 && b =1
\end{align*}
Veamos qué tiempos empíricos registramos para este algoritmo:

\begin{table}[H]
\renewcommand{\arraystretch}{1.2}
\centering
\begin{tabular}{|r|r|r|r|}
\hline
\rowcolor{LightSteelBlue2}
\multicolumn{1}{|c|}{n} & \multicolumn{1}{c|}{Tiempo (ms)} & \multicolumn{1}{c|}{Repeticiones} & \multicolumn{1}{c|}{$n^3/T_\text{medio}$} \\ \hline
100                     & 51                               & 10                                & 196.078,4314                                       \\ \hline
200                     & 384                              & 10                                & 208.333,3333                                       \\ \hline
400                     & 2.831                            & 10                                & 226.068,5270                                       \\ \hline
800                     & 22.475                           & 10                                & 227.808,6763                                       \\ \hline
1.600                   & 17.705                           & 1                                 & 231.347,0771                                       \\ \hline
3.200                   & 142.593                          & 1                                 & 229.800,9019                                       \\ \hline
6.400                & FdT                              & -                                 & -                                                  \\ \hline
\end{tabular}
\end{table}
Podemos observar que los tiempos empíricos respaldan la complejidad teórica.

\vspace{6mm}  
\result{Toma de tiempos para \texttt{sustraccion5}}
Análogamente, para este método:
\begin{table}[H]
\renewcommand{\arraystretch}{1.2}
\centering
\begin{tabular}{|r|r|r|r|}
\rowcolor{LightSteelBlue2}
\hline
\multicolumn{1}{|c|}{n} & \multicolumn{1}{c|}{Tiempo (ms)} & \multicolumn{1}{c|}{Repeticiones} & \multicolumn{1}{c|}{3\textasciicircum{}\{n/2\}/T} \\ \hline
30                      & 422                              & 1                                 & 34.002,1493                                       \\ \hline
32                      & 1.280                             & 1                                 & 33.630,2508                                       \\ \hline
34                      & 3.792                             & 1                                 & 34.055,9502                                       \\ \hline
36                      & 11.445                            & 1                                 & 33.850,6325                                       \\ \hline
38                      & 34.310                            & 1                                 & 33.875,2978                                       \\ \hline
40                      & 103.475                           & 1                                 & 33.696,8775                                       \\ \hline
42                      & FdT                              & -                                 & -                                                 \\ \hline
\end{tabular}
\end{table}
Sigue la complejidad esperada

\newpage
\result{Tiempo para \texttt{sustraccion5} con \texorpdfstring{$n=80$}{n=80}}
Razonando como en el \hyperref[result:1.3.]{apartado 1.3}, al comparar con la tabla anterior,
$$t_2 = 103.475ms \cdot \frac{3^{\frac{80}{2}}}{3^{\frac{40}{2}}}  \cdot\frac{1h}{3,6 \cdot 10^6}ms \cdot \frac{1\text{año}}{24*365h}= 1,144 \cdot 10^4 \text{años}$$

\newpage
\chapter{Divide y vencerás por división}
\result{Complejidades teóricas}

\end{document} 